这个路线定下来如何？
1. 完成自研 DUT 的基础验证闭环
先完成 Fable 独立审查、读写混合测试、在途复位、参数扫描、覆盖基线和关键故障注入。这一步交付的是一个结果可信、可以固定下来的基准版本。
同时整理验证环境的依赖：哪些是 AXI 通用规则，哪些是自研 Crossbar 的地址映射、容量或调度策略。
2. 把验证环境抽离为可复用组件，接入一个第三方 DUT
这是你的首要差异化目标，优先于增加 QoS 等新功能。

* Agent、接口协议检查尽量通用。
* Scoreboard、覆盖模型通过配置描述拓扑、地址映射、ID 和功能范围。
* 自研及第三方 DUT 各自使用薄封装接入相同测试环境。
* 先选一个第三方实现，跑双方共同支持的功能，再扩展。

这里的价值不仅是“外部 DUT 通过了测试”，还包括：**移植时发现了哪些自研假设、如何消除验证环境与 DUT 的耦合。**不必以发现开源 bug 为目标。
3. 围绕关键模块加入形式验证
先从 Tracker、FIFO、仲裁和 ID 路由入手，证明容量不越界、不凭空完成、背压稳定、ID 恢复和顺序保护等具体性质。
“无死锁”放到后面，并明确下游最终响应、背压最终解除等环境假设。先做几个边界清楚、能解释证明范围的结果，比一开始宣称证明整个 Crossbar 更有价值。
4. 建立统一流量模型，做定量性能研究
在同一套流量和统计口径下，比较自研与第三方实现，并扫描：

* 均匀访问与热点 Slave；
* outstanding 深度；
* burst 长度与读写比例；
* 下游背压；
* 仲裁等待、吞吐和平均／尾部延迟。

**先研究现有轮询仲裁的行为，再决定是否增加第二种策略。**这样每个新增设计都有数据依据。
5. 做综合和时序评估，选择一个有动机的扩展
比较时统一功能范围、参数、工艺库和约束；面积也要明确是逻辑资源估计还是映射后的单元面积。
扩展根据前面的结果选择：

* 关键路径限制频率 → 可选寄存器切片；
* 热点竞争导致延迟问题 → 仲裁／QoS；
* 想突出隔离与系统可靠性 → 地址权限或此前讨论的 Slave 故障隔离。

先做一个，给出功能收益、性能变化和资源代价，再考虑第二个。
6. 整理可复现交付和面试叙事
最终形成三条能深入展开的故事：

* **设计故事：**为什么这样组织事务跟踪、顺序保护和缓冲。
* **验证故事：**怎样发现 RTL／验证环境问题，如何通过第三方 DUT、故障注入和形式验证提高可信度。
* **取舍故事：**哪种负载暴露瓶颈，怎样优化，付出了什么代价。

考虑你偏向 verification 的目标，我会把优先级定为：
基础验证闭环 → 可复用 VIP＋一个第三方 DUT → 关键模块形式验证 → 性能与 PPA 对比 → 一个数据驱动的扩展。
这条路线保留完整功能目标，同时让项目的吸引力集中在“验证能力可信、设计取舍有证据”上。